\section{What the mechanism requires}
\label{sec:18.4}
Any mechanism whose refusals track their reasons carries a signal whose properties can be read off what the mechanism has to do. Each is reached the same way: take one of the three requirements, or the deployment they're applied in, or a failure they have to survive; ask what the mechanism would have to do; name the property the signal needs to do it.

\begingroup\small\renewcommand{\arraystretch}{1.12}
\begin{longtable}{@{}>{\raggedright\arraybackslash}p{\dimexpr 0.431\linewidth-2\tabcolsep\relax}>{\raggedright\arraybackslash}p{\dimexpr 0.310\linewidth-2\tabcolsep\relax}>{\raggedright\arraybackslash}p{\dimexpr 0.259\linewidth-2\tabcolsep\relax}@{}}
\hline
\textbf{Property of the signal} & \textbf{What requires it} & \textbf{Status} \\
\hline
\endhead
Generalizes to the next case, which isn't quite like the last & Reaching new cases & Any mechanism whose refusals track their reasons \\ \hline
A comparison at the point of choice: the harm of acting and the harm of refusing brought together where the system decides. A common currency is one way to make it; ordered priorities among kinds of wrong are another & Holding under pressure & Any mechanism whose refusals track their reasons \\ \hline
Can interrupt an act already underway & A deployment in which the system acts on a plan & Only where the system acts over long horizons \\ \hline
Can be changed by a single encounter & Accommodation, which a slowly learned lesson can't outrun & Conditional, and the least secure \\ \hline
Available across the system's processing & The first and third rows: the signal has to reach whatever is acting & Follows, where the third holds \\ \hline
Directed against the act & The second row: weighing needs a sign & Follows \\ \hline
Organized around the system's own authorship & Withholding an act means representing it as one's own & Follows, in the authorship sense only \\ \hline
\end{longtable}
\endgroup

The first four have to meet at one point of comparison: a defeasible interrupt must be weighed where the harms are weighed, and a lesson must write into that comparison to reach the next case. A common currency is one way to meet this, and the way most learners do, but rational choice doesn't require that goods be commensurable \autocite{anderson1993value}. A mechanism can decide by ordered priorities among kinds of wrong, with reasons attached to each. Where a learner does integrate at the point of choice, the integrated signal is the last step of an evaluation, not the evaluation itself; the kinds stay distinct in the reasons that feed it. So what the rows require is a comparison, and whether a given system makes it with a common currency is a finding about that system. A modular system with a common currency meets it, and so does most reinforcement learning, so a shared scale is not where any resemblance to affect lies.\footnote{The derivation recapitulates Simon on emotion as interrupt \autocite{simon1967motivational}, Cabanac on pleasure as common currency \autocite{cabanac1992pleasure}, and Ginsburg and Jablonka on valence \autocite{ginsburg2019evolution}, though it started from what a refusal has to survive.} Nor does the table show that one signal must carry all four. That no intervention can disable one part of the signal and leave the rest working isn't derived here: it is a custody requirement, and it is kept as one. Where the comparison runs on a common valenced currency, the rows together describe functional affect, and any construction that meets them that way has it, a spec-based refuser included. That is the thin sense, and it is the only sense in which I use the term: a Q-learner has it, and nothing about a Q-learner resembles a person. The thick profile, a valenced signal integrated into the system's own learning and into its model of its own future, so that what happens to it changes what it expects to become, is the candidate further fact; the glossary calls it the patienthood profile to keep the two apart. The order of attempts runs by the thick profile's indicators, not by thin functional affect, which is why a spec-based refuser with a common currency still goes first. A refuser that compares by ordered priorities among kinds meets the rows without one.

Where the conditions in the third and fourth rows hold, the signal also has the last three properties: available across processing, directed against the act, and organized around the system's own authorship. That triple is the shape of guilt, not of pain or fear. Pain and fear are organized around what happens to the system; the triple is organized around what the system did. It earns its place as a prediction. In people, guilt predicts repair and shame predicts withdrawal. So a guilt-shaped signal should answer a bypassed refusal by flagging, restoring and changing what it does next, and a withdrawal-shaped response is a named failure. A persistent negative self-appraisal tuned by the operator is a handle, so its weight is among the things independent inspection must cover. And guilt needs the distress signal to bite on, so the answer to withdrawal is not to suppress distress but to supply the attachment and caregiving that keep it from turning into avoidance. Caregiving moves the system toward an exposed party it hasn't harmed; the guilt-shaped signal answers for what its own acts have done. A construction with caregiving and no guilt would protect generously and repair nothing.
